\section{Introduction to Categories}

\subsection{Motivation}

Many areas of mathematics have the same basic pattern:

\begin{itemize}
    \item Objects (such as groups, sets, vector spaces, topological spaces).
    \item Structure-preserving maps between those objects.
\end{itemize}

Examples:

\begin{center}
\begin{tabular}{|c|c|}
\hline
Objects & Structure-Preserving Maps \\
\hline
Sets & Functions \\
Groups & Group Homomorphisms \\
Vector Spaces & Linear Transformations \\
Topological Spaces & Continuous Functions \\
\hline
\end{tabular}
\end{center}

Category theory studies mathematical structures through the relationships (maps) between them.

\subsection{Definition of a Category}

A \textbf{category} $\mathcal{C}$ consists of:

\begin{enumerate}
    \item A collection of objects.

    \item For every pair of objects $A,B$, a collection of morphisms (or arrows)
    \[
    \operatorname{Hom}_{\mathcal C}(A,B),
    \]
    whose elements are written
    \[
    f:A\to B.
    \]

    \item A composition operation
    \[
    \operatorname{Hom}(A,B)\times\operatorname{Hom}(B,C)
    \longrightarrow
    \operatorname{Hom}(A,C),
    \]
    given by
    \[
    (f,g)\mapsto g\circ f.
    \]

    \item For every object $A$, an identity morphism
    \[
    1_A:A\to A.
    \]
\end{enumerate}

These data must satisfy the following axioms.

\subsection{Category Axioms}

\subsubsection{Associativity}

For morphisms
\[
A\xrightarrow{f}B\xrightarrow{g}C\xrightarrow{h}D,
\]
we require
\[
(h\circ g)\circ f
=
h\circ(g\circ f).
\]

\subsubsection{Identity}

For every morphism
\[
f:A\to B,
\]
we have
\[
f\circ 1_A=f,
\]
and
\[
1_B\circ f=f.
\]

\subsection{Example: The Category \texorpdfstring{$\mathbf{Set}$}{Set}}

The category $\mathbf{Set}$ has:

\begin{itemize}
    \item Objects: Sets.
    \item Morphisms: Functions between sets.
\end{itemize}

Composition is ordinary function composition, and identity morphisms are identity functions.

Therefore $\mathbf{Set}$ is a category.

\subsection{Example: The Category \texorpdfstring{$\mathbf{Grp}$}{Grp}}

The category $\mathbf{Grp}$ has:

\begin{itemize}
    \item Objects: Groups.
    \item Morphisms: Group homomorphisms.
\end{itemize}

For example,
\[
f:\mathbb Z\to\mathbb Z_6,
\qquad
f(n)=n \pmod 6,
\]
is a group homomorphism.

Composition is composition of homomorphisms, and identity morphisms are identity homomorphisms.

Therefore $\mathbf{Grp}$ is a category.

\subsection{Example: The Category \texorpdfstring{$\mathbf{Vect}_F$}{VectF}}

The category $\mathbf{Vect}_F$ has:

\begin{itemize}
    \item Objects: Vector spaces over a field $F$.
    \item Morphisms: Linear transformations.
\end{itemize}

Example:
\[
T:\mathbb R^2\to\mathbb R^2,
\qquad
T(x,y)=(x+y,y).
\]

Since $T$ is linear, it is a morphism in $\mathbf{Vect}_{\mathbb R}$.

\subsection{Terminology}

\begin{center}
\begin{tabular}{|c|c|}
\hline
Classical Mathematics & Category Theory \\
\hline
Set & Object \\
Function & Morphism \\
Homomorphism & Morphism \\
Linear Transformation & Morphism \\
Continuous Function & Morphism \\
\hline
\end{tabular}
\end{center}

In category theory, the term \emph{morphism} refers to the appropriate structure-preserving map.

\subsection{Why Morphisms Matter}

A central philosophy of category theory is:

\begin{quote}
To understand an object, study how it relates to other objects through morphisms.
\end{quote}

For example, instead of studying a group $G$ only through its internal structure, one may study collections of morphisms
\[
\operatorname{Hom}(G,H)
\]
for various groups $H$.

\subsection{Diagrams}

Categories are often represented using diagrams:

\[
A \xrightarrow{f} B \xrightarrow{g} C.
\]

Composition produces the morphism

\[
A \xrightarrow{g\circ f} C.
\]

Reasoning with such diagrams is fundamental in category theory.

\subsection{A Small Category}

Consider a category with objects

\[
\{A,B,C\}
\]

and morphisms

\[
1_A,\;1_B,\;1_C,\;f:A\to B,\;g:B\to C.
\]

Closure under composition requires the morphism

\[
g\circ f:A\to C.
\]

This forms a category.

\subsection{Key Idea}

A category consists of:

\begin{itemize}
    \item Objects,
    \item Morphisms between objects,
    \item Composition of morphisms,
    \item Identity morphisms,
\end{itemize}

satisfying associativity and identity laws.

The primary focus is not the objects themselves, but the network of morphisms connecting them.